\chapter{Introducci\'on y objetivos del proyecto}\label{defobjetivos}

\section{Contexto y motivaci\'on}

El consumidor espa\~nol afronta la compra cotidiana en un contexto de fuerte variabilidad de
precios entre cadenas, ausencia de una fuente unificada y alto coste temporal para comparar
ofertas manualmente. Aunque el ahorro potencial entre supermercados es significativo, en la
pr\'actica ese ahorro se pierde cuando no se considera el coste de desplazamiento ni el tiempo de
ruta.

BarGAIN nace para resolver ese problema como una plataforma de compra inteligente que integra
precio, distancia y tiempo en una \'unica decisi\'on. El sistema combina una app m\'ovil para
consumidores, un portal web para peque\~nas y medianas empresas (PYME), y un backend modular
con capacidades geoespaciales, reconocimiento \'optico de caracteres (OCR) y asistencia
conversacional.

\section{Problema abordado}

Las soluciones actuales de comparaci\'on de supermercados suelen responder a la pregunta
``qu\'e tienda tiene el precio m\'as bajo'', pero no a la pregunta completa ``qu\'e compra me
conviene realmente seg\'un mi ubicaci\'on y mis preferencias''.

El problema de BarGAIN se formula como una optimizaci\'on multicriterio en la que intervienen:

\begin{itemize}
	\item coste total de la cesta,
	\item distancia adicional de desplazamiento,
	\item tiempo total estimado de compra y ruta.
\end{itemize}

\section{Objetivo principal}

Desarrollar una aplicaci\'on m\'ovil y web que optimice la cesta de la compra del usuario,
calculando la combinaci\'on de supermercados y comercios locales que ofrece la mejor relaci\'on
precio--distancia--tiempo seg\'un preferencias configurables.

\section{Objetivos espec\'ificos}

\begin{enumerate}
	\item Ofrecer un servicio centralizado que re\'una productos, tiendas, precios y listas de la
		compra, accesible tanto desde la aplicaci\'on de consumo como desde el portal de comercios.
	\item Localizar las tiendas pr\'oximas al usuario y preparar el c\'alculo de rutas a partir de
		su ubicaci\'on.
	\item Recomendar la mejor combinaci\'on de tiendas ponderando precio, distancia y tiempo seg\'un
		las preferencias que cada usuario configure.
	\item Permitir crear la lista de la compra a partir de una foto de un ticket o de una nota
		manuscrita, con revisi\'on y correcci\'on por parte del usuario antes de guardarla.
	\item Proporcionar un asistente conversacional que responda en lenguaje natural a consultas
		relacionadas con la compra, ce\~nido a ese dominio.
	\item Integrar a las grandes cadenas de supermercados en la comparativa manteniendo su
		cat\'alogo de precios actualizado de forma autom\'atica y peri\'odica, de modo que las
		recomendaciones partan siempre de datos reales y vigentes.
	\item Ofrecer al peque\~no comercio un portal propio desde el que gestionar sus precios y
		promociones.
	\item Garantizar la calidad mediante pruebas automatizadas a distintos niveles y ensayos
		manuales en los flujos que dependen del dispositivo m\'ovil.
	\item Disponer de un despliegue p\'ublico y reproducible que permita probar el sistema en
		condiciones reales.
\end{enumerate}

\section{Alcance del proyecto}

El alcance del Trabajo Fin de Grado (TFG) cubre el ciclo completo de ingenier\'ia del software:

\begin{itemize}
	\item an\'alisis y requisitos,
	\item dise\~no arquitect\'onico y de datos,
	\item implementaci\'on backend y frontend,
	\item pruebas funcionales y no funcionales,
	\item despliegue en entorno staging,
	\item documentaci\'on t\'ecnica y memoria final.
\end{itemize}

Quedan fuera del alcance de esta versi\'on: despliegue productivo global,
acuerdos formales con todas las cadenas para interfaces de programaci\'on de aplicaciones (API)
oficiales, y la versi\'on final de Live Activities nativas en iOS (registro de decisi\'on
arquitect\'onica, ADR-010).

\section{Tecnolog\'ias clave}

El sistema se apoya en un stack full-stack moderno y coherente:

\begin{itemize}
	\item \textbf{Backend:} Django 5 con Django REST Framework (DRF), Celery y Redis.
	\item \textbf{Datos:} PostgreSQL 16 + PostGIS 3.4.
	\item \textbf{Frontend:} React Native con Expo + companion web en React.
	\item \textbf{Inteligencia artificial (IA):} Google Cloud Vision API (OCR) y Gemini v\'ia
		backend proxy.
	\item \textbf{Optimizaci\'on:} OR-Tools + c\'alculo de distancias (OpenRouteService +
		fallback haversine).
	\item \textbf{Ops:} Docker, GitHub Actions, Render, Sentry.
\end{itemize}

\section{Estructura de la memoria}

La memoria se organiza para reflejar el recorrido completo del proyecto: contexto y objetivos,
estado del arte, comparativa, planificaci\'on, requisitos, dise\~no e implementaci\'on,
manual de usuario, pruebas y resultados, y conclusiones. Cierran el documento los ap\'endices
---con la relaci\'on de endpoints de la API y las capturas de la versi\'on de escritorio--- y las
referencias bibliogr\'aficas. Al inicio se incluye un glosario con los acr\'onimos empleados.
